\section{Introduction}

LLMs are now widely used in software engineering, particularly to assist developers with tasks such as program understanding, debugging, code review, and code generation. To evaluate their capabilities on these tasks, many benchmarks have been proposed with the objective of reproducing problems that are close to those encountered by developers.

However, these benchmarks mainly provide a global measure of whether a model can solve a specific task. They can therefore indicate whether a model performs well in a given scenario, but they make it more difficult to diagnose the reasons behind its success or failure. A single task can involve several skills at the same time, which makes it difficult to identify precisely which elements cause difficulties for the model. For example, a model may understand a program by partially relying on semantic information present in method or variable names, while its performance may decrease when this information is removed. A correct or incorrect answer alone is not sufficient to highlight this type of behavior.

In this work, we propose \benchmark, a Java benchmark whose objective is not to replace existing benchmarks, but to provide an evaluation framework that allows a more fine-grained analysis of model capabilities. \benchmark focuses on reachability problems between methods. Given a source method and a target method, the model must determine whether the target can be reached from the source through a sequence of calls.

The Java code used in the benchmark is automatically generated from call graphs. This provides an executable-graph oracle and the expected answer for each question. The model returns a JSON object with a Boolean final answer and an ordered sequence of explored methods. The evaluator scores these outputs independently, accepting any path order justified by executable edges rather than requiring one canonical traversal.

\benchmark provides two families of tasks. The first, \emph{unguarded reachability}, considers programs in which calls are executed whenever they are reached because they are outside control structures. The second, \emph{condition-aware reachability}, introduces calls inside conditions that can make some calls executable or non-executable. The model must then reason about the data flow to determine which conditions are satisfied and, as a consequence, which method calls are actually executed. This second task therefore distinguishes what is visible in the source code from what happens during program execution.

Through this approach, \benchmark aims to provide a more controlled and diagnostic evaluation of LLMs on program-understanding tasks. It makes it possible to study model behavior according to specific factors rather than limiting the evaluation to whether the final answer is correct.
This paper makes three contributions:
\begin{itemize}[leftmargin=*]
    \item a configurable benchmark that generates Java reachability problems together with executable call graphs and oracle labels;
    \item an empirical study of \texttt{[N\_MODELS]}, supported by an artifact containing the complete generation and evaluation pipeline.
\end{itemize}
